\begin{frame}[allowframebreaks]{}
    \LARGE Normalizing Flow Models: \\[1.5ex] \textbf{Limitations}
\end{frame}

\begin{frame}[allowframebreaks]{Limitations}
    \begin{itemize}
        \item \textbf{Invertibility constrains the architecture}: every layer must be invertible, which rules out most standard building blocks.
        \item \textbf{The determinant must stay cheap}: coupling layers make it $O(n)$, but only by restricting each layer to transform half its input at a time, so many layers are needed for expressiveness. These first two constraints apply to discrete-layer flows; flow matching drops both.
        \item \textbf{No dimensionality reduction}: the change of variables formula needs $f:\mathbb{R}^n \rightarrow \mathbb{R}^n$, so the latent space is exactly as large as the data. Unlike a VAE, a flow has no bottleneck and cannot compress.
        \item \textbf{Continuous data only}: the formula assumes continuous densities, so discrete values such as 8-bit pixels must first be dequantized by adding noise. RealNVP then applies a logit transform, $\mathrm{logit}(\alpha + (1-\alpha)x/256)$ with $\alpha = 0.05$, so that the density is not pushed against the boundaries of the pixel range.
        \item In practice flows can trail autoregressive models on likelihood, and large ones are slow.
    \end{itemize}
\end{frame}